\chapter{Implementation}
\label{ch:implementation}
\begin{tcolorbox}
[colback=magenta!5!white,colframe=magenta!75!black,
  title=About this chapter]
  This chapter details the implementation of the Type3 Asset Administration Shell (AAS) integrated with a Multi-Agent System (MAS) using the Belief-Desire-Intention (BDI) approach. It covers the technical aspects of the AAS development, including the software architecture, key components, and the integration process with the MAS. The chapter also discusses the challenges encountered during implementation and the solutions adopted to address them.
\end{tcolorbox}
\section{AAS implementation and repository design pattern}
\begin{tcolorbox}
[colback=magenta!5!white,colframe=magenta!75!black,
  title=About AAS Implementation and Repository Design Pattern]
  This section describes the implementation of the Type3 Asset Administration Shell (AAS) and the application of the repository design pattern. It outlines the software architecture, key components, and data management strategies employed in the AAS development. The section also discusses how the repository design pattern facilitates efficient data access, storage, and retrieval within the AAS framework, enhancing its scalability and maintainability.
\end{tcolorbox}
\section{Resource Agents implementation}
\begin{tcolorbox}
[colback=magenta!5!white,colframe=magenta!75!black,
  title=About Resource Agents Implementation]
  This section details the implementation of Resource Agents within the Multi-Agent System (MAS) integrated with
  the Type3 Asset Administration Shell (AAS). It discusses the design principles, architectural considerations, and interaction mechanisms that enable Resource Agents to effectively manage and represent physical assets. The section also highlights the role of Resource Agents in facilitating communication and cooperation among agents, as well as their contribution to the overall functionality and autonomy of the integrated AAS-MAS framework.
\end{tcolorbox}
\section{Production Agents implementation}
\begin{tcolorbox}
[colback=magenta!5!white,colframe=magenta!75!black,
  title=About Production Agents Implementation]
  This section outlines the implementation of Production Agents within the Multi-Agent System (MAS) integrated with
  the Type3 Asset Administration Shell (AAS). It discusses the design principles, architectural considerations, and interaction mechanisms that enable Production Agents to effectively manage and represent product-related information. The section also
  highlights the role of Production Agents in facilitating communication and cooperation among agents, as well as their contribution to the overall functionality and autonomy of the integrated AAS-MAS framework.
\end{tcolorbox}
\section{MAS communication over XMPP protocols in Asynchronous way}
\begin{tcolorbox}
[colback=magenta!5!white,colframe=magenta!75!black,
  title=About MAS Communication over XMPP Protocols in Asynchronous Way]
  This section discusses the implementation of communication mechanisms within the Multi-Agent System (MAS) integrated with
  the Type3 Asset Administration Shell (AAS) using XMPP protocols in an asynchronous manner
. It outlines the architectural design, communication protocols, and interaction patterns that facilitate efficient and reliable message exchange among agents. The section also highlights the benefits of using XMPP for asynchronous communication, including scalability, flexibility, and real-time data transfer, which enhance the overall performance and responsiveness of the integrated AAS-MAS framework.
\end{tcolorbox}